\cotefile{157-5}
\scanpage{421}
\mspage{333}

\[
\begin{aligned}
x&=\varinjlim_{i\in I}x_i,\qquad \text{les }x_i\text{ dans }\underline S(\pi)\\
&& I\ \unsure{\text{grd}|\text{gd}|\text{gros}}\text{ devant }\pi.
\end{aligned}
\]

Alors le foncteur composé $i\mapsto x_i$
\diagram{\begin{tikzcd}[ampersand replacement=\&]
 I \arrow[r] \& \underline S(\pi)_{/x} \arrow[r,hook] \& \unsure{M|\mathcal M}_{\pi/x}
\end{tikzcd}}
est cofinal\inserted{(cf. 4.4.2 \unsure{pag.|page|p.} 79)}, \deleted{et comme} \inserted{donc} $\underline S(\pi)_{/x}$ \unsure{$-\!-$|$\longrightarrow$} $M_{\pi/x}$ \deleted{est \unsure{pleint|pleint\textsuperscript{t}|plein\textsuperscript{t}} fidèle, on voit qu'il} satisfait la condition F1 (puisque le composé la satisfait), et comme $\underline S(\pi)_{/x}\longrightarrow M_{\pi/x}$ est pl. fid, cela implique déjà que ce foncteur est cofinal. Dans \deleted{\unsure{par}} \unsure{cas}, on en conclut $x\in\mathrm{Ob}\,\underline S$, qed.

\deleted{\unsure{D}}Supposons maintenant que $S$ satisfasse \unsure{a')|d')} (au lieu de b)), et prouvons encore l'équivalence \unsure{$(*)$|$(**)$}. Pour ceci, écrivons
\[
 x=\varinjlim_{j\in\unsure{J|I}}x_j\qquad
 \left|\begin{aligned}
 &\text{les }x_j\in\mathrm{Ob}\,M_\pi\\
 &J\ \text{\uline{ens. ordonné}}\quad\unsure{*}\\
 &(x_{\unsure{i|j}})_{\unsure{i|j}\in\unsure{J|I}}\ \deleted{\text{est}}\ \text{est }\unsure{\pi_0|\text{bon}} \text{-}\pi\text{-adapté}\\[-2pt]
 &\deleted{\unsure{=|-}\,\unsure{\text{grd}|\text{gd}}\text{ devant }\pi}
 \end{aligned}\right.
\]
% slanted
\margin{\uline{NB} Ça existe, \unsure{grâce|après|appliq.} \unsure{à|?} 4.13.14. $\square$}
ce qui est possible puisque \unsure{$\pi$|il} est \unsure{utile}, pour $M$. \unsure{$]$} Définissons $J_{\underline S}$ comme dans l'énoncé de a), donc comme image inverse de $\underline S(\pi)_{/x}$ $\unsure{\subset|\hookrightarrow}\,M_{\pi/x}$ par $I\unsure{\longrightarrow|-}M_\pi$. Par hyp. \unsure{d')|a')}, on a
\[
x\unsure{\mathord{+}|\in}\mathrm{Ob}\,\underline S\quad\Longleftrightarrow\quad I_{\underline S}\text{ est cofinal dans }I
\tag*{$(**')$}
\]

En utilisant $(**')$, prouvons \inserted{\unsure{encore|alors}} $(*)$\inserted{(p. \deleted{\unsure{0}}\unsure{331})}. \deleted{\unsure{P}} En fait, on utilisera \unsure{seulement} qu'on\unsure{.| a}

% archive sheet number: 420

